\section{Модуль I. Векторные пространства и евклидова геометрия}

\subsection{Практический вопрос I.1: Как проверить, что subset является subspace?}
\textbf{Ответ:} Для проверки, является ли подмножество $U \subseteq V$ подпространством, применяется стандартный
трехшаговый алгоритм:
\begin{enumerate}
    \item \textbf{Проверка нуля:} Проверить, содержит ли $U$ нулевой вектор ($0 \in U$).
    Если $0 \notin U$, множество гарантированно не является подпространством.
    \item \textbf{Замкнутость по сложению:} Взять два произвольных элемента $u_1, u_2 \in U$ и проверить, что их
    сумма $u_1 + u_2 \in U$.
    \item \textbf{Замкнутость по умножению на скаляр:} Взять произвольный $u \in U$ и скаляр $c \in \mathbb{R}$
    и проверить, что $c u \in U$.
\end{enumerate}
\emph{Объединенный критерий:} $\forall u_1, u_2 \in U, \forall \alpha, \beta \in \mathbb{R} \implies \alpha u_1 + \beta
u_2 \in U$.

\subsection{Практический вопрос I.2: Как span, independence и basis работают в вычислениях?}
\textbf{Ответ:} В вычислительной практике все три концепции реализуются через метод Гаусса над матрицей $A$,
составленной из исследуемых векторов как из столбцов:
\begin{enumerate}
    \item \textbf{Проверка линейной независимости:} Приводим $A$ к ступенчатому виду.
    Если все столбцы содержат ведущие элементы (нет свободных переменных), система независима.
    \item \textbf{Нахождение базиса $\spanop(v_1, \dots, v_k)$:} Столбцы исходной матрицы $A$, соответствующие
    ведущим позициям ступенчатого вида, образуют базис оболочки.
    \item \textbf{Проверка принадлежности вектора $b \in \spanop(v_1, \dots, v_k)$:} Решаем систему $Ax = b$.
    Если система совместна , $b \in \spanop$; коэффициенты $x$ дают координаты разложения.
\end{enumerate}

\subsection{Практический вопрос I.3: Как меняются coordinates при смене basis?}
\textbf{Ответ:} Пусть вектор $x$ имеет координаты $[x]_B$ в базисе $B$ и $[x]_C$ в базисе $C$.
Строим матрицу перехода $P_{C \leftarrow B}$, столбцами которой являются координаты векторов старого базиса $B$,
выраженные в новом базисе $C$:
\[ [x]_C = P_{C \leftarrow B} [x]_B. \]
Обратный пересчет осуществляется с помощью обратной матрицы: $[x]_B = (P_{C \leftarrow B})^{-1} [x]_C$.

\subsection{Практический вопрос I.4: Как соотносятся dot product, norm, distance и angle?}
\textbf{Ответ:} Скалярное произведение $\langle x, y \rangle$ является первоисточником всех метрических величин:
\begin{itemize}
    \item \textbf{Норма (длина):} $\|x\| = \sqrt{\langle x, x \rangle}$;
    \item \textbf{Расстояние:} $d(x, y) = \|x - y\| = \sqrt{\langle x-y, x-y \rangle}$;
    \item \textbf{Угол:} $\cos \theta = \frac{\langle x, y \rangle}{\|x\| \|y\|}$ (корректность формулы
    $|\cos \theta| \le 1$ обеспечена неравенством Коши-Буняковского-Шварца);
    \item \textbf{Ортогональность:} $\theta = 90^\circ \iff \langle x, y \rangle = 0$.
\end{itemize}

\subsection{Практический вопрос I.5: Как работает Principal component analysis (PCA)?}
\textbf{Ответ:} Пошаговый алгоритм PCA:
\begin{enumerate}
    \item \textbf{Центрирование данных:} Вычесть среднее по каждому признаку: $X_c = X - \mu$.
    \item \textbf{Ковариационная матрица:} Вычислить $\Sigma=\frac{1}{n-1} X_c^\top X_c \in \mathbb{R}^{d \times d}$
    \item \textbf{Спектральное разложение / SVD:} Найти собственные векторы $v_i$ матрицы $\Sigma$
    (или правые сингулярные векторы $X_c$).
    \item \textbf{Сортировка:} Упорядочить компоненты по убыванию собственных значений $\lambda_1 \ge \lambda_2
    \dots \ge \lambda_d$.
    \item \textbf{Проекция:} Выбрать $k$ первых компонент $V_k = [v_1 \dots v_k]$ и вычислить сжатые координаты
    $Z = X_c V_k \in \mathbb{R}^{n \times k}$.
\end{enumerate}

\newpage

\section{Модуль II. Матрицы, линейные отображения, ранг и СЛУ}

\subsection{Практический вопрос II.1: Как построить matrix of a linear transformation?}
\textbf{Ответ:} Алгоритм построения матрицы отображения $T \colon V \to W$:
\begin{enumerate}
    \item Зафиксировать базис $\mathcal{B} = (v_1,\dots,v_n)$ в $V$ и базис $\mathcal{C} = (w_1, \dots, w_m)$ в $W$.
    \item Подействовать оператором $T$ на каждый базисный вектор $v_j$ ($j = 1, \dots, n$).
    \item Разложить полученный вектор $T(v_j)$ по базису $\mathcal{C}$: $T(v_j) = \sum_{i=1}^m a_{ij} w_i$.
    \item Записать найденные скаляры $a_{1j}, \dots, a_{mj}$ в качестве $j$-го столбца результирующей матрицы.
\end{enumerate}

\subsection{Практический вопрос II.2: Как determinant, rank и inverse диагностируют одну проблему?}
\textbf{Ответ:} Все три величины диагностируют проблему потери информации (вырожденности):
\begin{itemize}
    \item $\det A = 0 \iff$ объём сжимается в нуль;
    \item $\rank A < n \iff$ образ имеет размерность меньше размерности прообраза;
    \item Несуществование $A^{-1} \iff$ преобразование невозможно обратить вспять из-за схлопывания ядра в ноль.
\end{itemize}

\subsection{Практический вопрос II.3: Как определить число решений $Ax = b$?}
\textbf{Ответ:} По теореме Кронекера-Капелли приводим расширенную матрицу $(A \mid b)$ к ступенчатому виду:
\begin{enumerate}
    \item Если $\rank(A) < \rank(A \mid b)$ — \textbf{нет решений} (система несовместна, возникла строка вида
    $(0 \dots 0 \mid c)$ при $c \ne 0$).
    \item Если $\rank(A) = \rank(A \mid b) = n$ (числу неизвестных) — \textbf{ровно одно решение}.
    \item Если $\rank(A) = \rank(A \mid b) = r < n$ — \textbf{бесконечно много решений}
    (размерность аффинного пространства решений равна $n - r$).
\end{enumerate}

\subsection{Практический вопрос II.4: Почему LU полезно для systems?}
\textbf{Ответ:} Факторизация $A = LU$ требует $\mathcal{O}\left(\frac{2}{3} n^3\right)$ операций и выполняется один раз.
После этого решение системы $Ax = b$ для любой новой правой части $b$ сводится к решению двух треугольных систем ($Ly =
b$ и $Ux = y$), каждая из которых требует лишь $\mathcal{O}(n^2)$ операций.
Прямой метод Гаусса требовал бы $\mathcal{O}(n^3)$ на каждый вектор $b$.

\subsection{Практический вопрос II.5: Какие common dimensional mistakes возникают с matrices?}
\textbf{Ответ:}
\begin{enumerate}
    \item Умножение $A_{(m \times k)} B_{(p \times n)}$ при $k \ne p$.
    \item Сложение матриц разных размеров.
    \item Ошибочное транспонирование: $(AB)^\top \ne A^\top B^\top$ (верно $B^\top A^\top$).
    \item Путаница между внутренним произведением $x^\top y\in\mathbb{R}$ и внешним
    $x y^\top \in\mathbb{R}^{n \times n}$.
\end{enumerate}

\newpage

\section{Модуль III. Собственные значения и матричные разложения}

\subsection{Практический вопрос III.1: Как compute eigenvectors after eigenvalues?}
\textbf{Ответ:}
\begin{enumerate}
    \item Подставить найденное значение $\lambda_i$ в однородную систему $(A - \lambda_i I) v = 0$.
    \item Привести матрицу $A - \lambda_i I$ методом Гаусса к ступенчатому виду.
    \item Разделить переменные на базисные (ведущие) и свободные.
    \item Выразить базисные переменные через свободные и записать фундаментальную систему решений (ФСР), столбцы
    которой образуют базис собственного подпространства $E_{\lambda_i}$.
\end{enumerate}

\subsection{Практический вопрос III.2: Как диагонализация упрощает возведение в степень?}
\textbf{Ответ:} Если $A = P \Lambda P^{-1}$, то:
\[ A^k = (P \Lambda P^{-1})^k = P \Lambda^k P^{-1} = P \diag(\lambda_1^k, \dots, \lambda_n^k) P^{-1}. \]
Вместо $k - 1$ матричных умножений со сложностью $\mathcal{O}(k n^3)$ вычисляются лишь степени чисел $\lambda_i^k$, что
требует $\mathcal{O}(n^3 + n \log k)$ операций.

\subsection{Практический вопрос III.3: Как геометрически интерпретировать eigenvectors?}
\textbf{Ответ:} Собственный вектор $v$ указывает направление в пространстве, вдоль которого преобразование действует как
\textbf{чистое одномерное растяжение или сжатие} на коэффициент $\lambda$.
Вектор $Av$ остаётся на той же прямой, что и вектор $v$ (прямая инвариантна относительно оператора).

\subsection{Практический вопрос III.4: Почему symmetric matrices удобны?}
\textbf{Ответ:} 1) Все собственные значения гарантированно вещественны.
2) Собственные векторы ортогональны.
3) Всегда существует ортонормированный базис из собственных векторов ($A = Q \Lambda Q^\top$).
4) Обращение матрицы $Q$ тривиально и численно устойчиво: $Q^{-1} = Q^\top$.

\subsection{Практический вопрос III.5: Чем decompositions отличаются с точки зрения применимости?}
\textbf{Ответ:}
\begin{itemize}
    \item $LU$: общего назначения для квадратных матриц при решении СЛУ.
    \item Cholesky ($LL^\top$): в 2 раза быстрее $LU$, для симметричных положительно определенных матриц
    (МНК, оптимизация).
    \item $QR$: численно сверхустойчив, для переопределенных систем и вычисления собственных значений.
    \item Спектральное: для анализа квадратичных форм, физических мод и динамики.
    \item $SVD$: универсально для любых прямоугольных матриц, оптимально для низкоранговой аппроксимации и сжатия.
\end{itemize}

\newpage

\section{Модуль IV. SVD и итоговые инструменты}

\subsection{Практический вопрос IV.1: Опишите способ вычисления матрицы SVD}
\textbf{Ответ:}
\begin{enumerate}
    \item Вычислить симметричную матрицу $A^\top A$.
    \item Найти ее собственные значения $\lambda_i \ge 0$ и упорядочить $\lambda_1 \ge \dots \ge \lambda_r > 0$.
    \item Сингулярные числа: $\sigma_i = \sqrt{\lambda_i}$.
    \item Найти ортонормированные собственные векторы $v_i$ матрицы $A^\top A$ $\implies$ матрица $V$.
    \item Вычислить левые сингулярные векторы: $u_i = \frac{1}{\sigma_i} A v_i$ для $i = 1, \dots, r$.
    \item При необходимости дополнить $\{u_i\}$ до ортонормированного базиса $\mathbb{R}^m \implies$ матрица $U$.
    \item Собрать $A = U \Sigma V^\top$.
\end{enumerate}

\subsection{Практический вопрос IV.2: Как reduced SVD используется в рекомендательных системах?}
\textbf{Ответ:} Матрица оценок пользователей $R_{m \times n}$ факторизуется через усеченное SVD: $R \approx U_k \Sigma_k
V_k^\top = P Q^\top$.
Строки матрицы $P = U_k \Sigma_k^{1/2}$ отражают вкусовые профили пользователей в $k$ латентных темах, а строки $Q = V_k
\Sigma_k^{1/2}$ — профили контента.
Пропущенная оценка восстанавливается как $\hat{r}_{ui} = \langle p_u, q_i \rangle$.

\subsection{Практический вопрос IV.3: Чем полезна low-rank approximation? Назовите пример}
\textbf{Ответ:} Низкоранговая аппроксимация сохраняет главную энергию матрицы, отсекая шум и радикально сокращая память
с $m \cdot n$ до $k(m + n + 1)$.
\emph{Пример:} Сжатие изображений: фотография $1000 \times 1000$ при $k = 50$ занимает в 10 раз меньше памяти при почти
неотличимом качестве.

\subsection{Практический вопрос IV.4: Что такое pseudoinverse matrix, чем она полезна?}
\textbf{Ответ:} Псевдообратная матрица Мура-Пенроуза $A^+ = V \Sigma^+ U^\top$ обобщает понятие обратной матрицы на
случай любых прямоугольных или вырожденных матриц.
Она даёт наилучшее обобщенное решение системы $Ax = b$: $x^* = A^+ b$.
Если решений нет, $x^*$ минимизирует ошибку $\|Ax - b\|_2$; если решений бесконечно много, $x^*$ выбирает решение с
наименьшей нормой $\|x\|_2$.

\newpage

% =========================================================================
% ЧАСТЬ 3: БАЗОВЫЕ ВОПРОСЫ (BASIC)
% =========================================================================